% page 149 of the facsimile (image p-148 of the scan)
% block 167.6 x 242.0 mm ; left margin 34.4 mm
\pgc{12.700mm}{0.000mm}{
\l{}
\l{}
\l{\cel{54}139}
\l{}
\l{}
\l{epigrafo. I. La skriburo quan onu renkontras o pozas sur la}
\l{\cel{3}monumenti. - II. Kurta cito-texto quan onu lokizas kape di}
\l{\cel{3}libro, di chapitro, por indikar lua tendenco. - DEFIRS.}
\l{}
\l{\sou{epigrafio}. La cienco di la epigrafi. - DEFIRS.}
\l{}
\l{\sou{epigramo}. Mikra versajo qua kontenas espritajo pikanta.-DEFIRS.}
\l{}
\l{\sou{epika}. Qua naracas per versi ago heroala. - DEFIRS.}
\l{}
\l{\sou{epikarpo}. Shelo extera di la frukto.}
\l{}
\l{\sou{epilepsio}. (\sou{patol.}) Morbo di la encefalo, karakterizata da acesi}
\l{\cel{3}dum qui la malado esvanas subite e kaptesas da konvulsi. -}
\l{\cel{3}DEFIRS.}
\l{}
\l{\sou{epilogo}.\cel{2}I. (\sou{che la antiqui}). Kurta diskurso recitita da la}
\l{\cel{3}aktoro lor la fino di la pleo, e per qua lu demandis de la}
\l{\cel{3}spektinti aplaudi. - II. Rezumo, konkluzo, ye la fino di poemo,}
\l{\cel{3}di diskurso, di libro. - DEFIRS.}
\l{}
\l{\sou{epistologismo}. (\sou{filoz.}) Sorito du-silogisma, rezono dum qua onu}
\l{\cel{3}prenas kom un ek la premisi lo konkluzita de silogismo pre-}
\l{\cel{3}iranta.}
\l{}
\l{\sou{episkopo}. Alt-oficiero di la Eklezio kristana, chefo e pastoro di}
\l{\cel{3}diocezo. - DEFIRS.}
\l{}
\l{\sou{epistemologio}. Kritikado di la cienci. - DEFIS.}
\l{}
\l{\sou{epistolo}. I. (\sou{literaturo}). Letro versa. - II. (\sou{liturgio katolika}).}
\l{\cel{3}Extraktajo de la letri kanonala o de la cetera libri santa,}
\l{\cel{3}quan onu lektas o kantas lor la meso, ante la evangelio. -}
\l{\cel{3}DEFIRS.}
\l{}
\l{\sou{epitafo}. Skriburo sur tombo. - DEFIRS.}
\l{}
\l{\sou{epitalamio}. Mikra poemo, kompozita okazione mariajo, honore di}
\l{\cel{3}la jusa spozi. - DEFIS.}
\l{}
\l{\sou{epitelio}. (\sou{anat.}) Quaza epidermo qua kovras la membrani mukoza,}
\l{\cel{3}seroza, vaskulala e glandala.- DEFIRS.}
\l{}
\l{\sou{epiteto}. (\sou{gram.}) Vorto quan onu adjuntas a substantivo por plu-}
\l{\cel{3}saliigar la ideo quan ica expresas. - DEFIRS.}
\l{}
\l{\sou{epitomo}. Abreviuro de libro. - DEFIRS.}
\l{}
\l{\sou{epizodo}. I. Agado acesora qua ne ligesas rigoroze a la temo.}
\l{\cel{3}- II. Fakto acesora qua ligesas plu o min rigoroze ad ensemblo}
\l{\cel{3}de fakti. - DEFIRS.}
}
